\documentclass[11pt,a4paper]{article}

\usepackage[utf8]{inputenc}
\usepackage[T1]{fontenc}
\usepackage{lmodern}
\usepackage{amsmath}
\usepackage{amssymb}
\usepackage{amsthm}
\usepackage[hidelinks]{hyperref}
\usepackage{doi}
\input{glyphtounicode}
\pdfgentounicode=1
\hypersetup{
  pdftitle={Spin as a Transformation Class of Rotating Wave Modes},
  pdfauthor={Jerome Beau},
  pdfkeywords={spin, helicity, spinor, SU(2) representations, univalence, superselection,
               neutron interferometry, rotating wave modes, Riemann-Silberstein vector,
               chiral transversality, composite Higgs}
}

\newcommand{\ket}[1]{\left|#1\right\rangle}

\newtheorem{proposition}{Proposition}
\newtheorem{lemma}{Lemma}
\newtheorem{corollary}{Corollary}

\title{Spin as a Transformation Class of Rotating Wave Modes}
\author{J\'er\^ome Beau}
\date{July 2026}

\begin{document}

\maketitle

\begin{abstract}
This note examines the hypothesis that the wave function is a combination of rotating and non-rotating wave
modes, spin expressing the rotation of these components rather than the rotation of a particle-like object.
For the photon the hypothesis is nearly literal: the transverse components of a circularly polarized wave
rotate, and the associated angular momentum is mechanically measurable.
A scalar counterexample shows that the rotation of components is not sufficient: spin is fixed by the way
the components transform under rotations of physical space.
The resulting hierarchy is: the relative phase is a measurable coordinate within a sector, the weight $m$ is
an axial winding degree, the spin $j$ labels the irreducible representation of $SU(2)$, and univalence
$(-1)^{2j}$ separates superselection sectors as a rule on the algebra of observables.
The formulation avoids the historical superluminal-velocity objection, which constrains rigid bodies rather
than internal rotations of wave components.
Any rotational ontology of spin must involve a relational, extended, or topologically attached wave
configuration, never a point particle endowed with a mere internal axis.
A soldering audit sharpens the programme-facing conclusion into two formal results: an exact
Veronese-type obstruction excludes the first internal candidate for selecting the compact real form of
the carrier square, and the currently derived data are shown to leave the spinorial soldering
indeterminate — the missing data being a phase-sensitive real-form selector, an independent rotation
action on the emergent space, and a Casimir-independent normalisation.
A separate, independent downstream reconnaissance shows that the Dirac-conjugated fermionic bilinear
$\overline{\psi_R} \psi_L$ carries exactly the Higgs doublet's electroweak quantum numbers in every
tree-level Yukawa sector, so that spinorial chiral structure permits, but does not by itself generate,
a composite Higgs condensate.
Interpretive outlook (not a result): these findings leave open two structural readings.
A future derived and sufficient soldering could make rotations effective symmetries of projected
geometry, whereas an irreducible residual indeterminacy would be consistent with treating the spinorial
carrier as prior to space.
\end{abstract}

\vspace{1em}
\noindent\textbf{Keywords:} spin; helicity; spinor; SU(2) representations; univalence; superselection;
neutron interferometry; rotating wave modes; Riemann--Silberstein vector; chiral transversality;
composite Higgs.

\clearpage
\tableofcontents

\clearpage
\section{Introduction}
\label{sec:introduction}

Spin is commonly introduced through negations: a rotation that is not a rotation, an angular momentum that
is not orbital, a magnetic moment without a current loop.
Faced with these negations, a natural hypothesis arises, and it is the founding hypothesis of this note:
the wave function is a combination of rotating and non-rotating wave modes, and spin expresses the rotation
of these components.
The rotating entity is from the outset a wave, not a particle-like object endowed with an internal axis;
the image is suggested by circularly polarized light, where two transverse field components in quadrature
produce a genuinely rotating field vector.
The task is then to determine which features of this rotation are physical, which are measurable, and what
the hypothesis must still explain.

This note assesses that hypothesis systematically.
Section~\ref{sec:photon} treats the photon, where the hypothesis is at its strongest.
Section~\ref{sec:rotatingwaves} states the hypothesis in its rigorous form and identifies the
transformation law of the components under spatial rotations as the decisive criterion.
Section~\ref{sec:representations} recalls why discrete spin values cannot be identified with particular
phase angles, and locates the correct discrete structure in the representation theory of $SU(2)$.
Section~\ref{sec:interferometry} presents the interferometric observability of the spinor sign.
Section~\ref{sec:univalence} formulates the univalence superselection rule as a statement about the algebra
of observables.
Section~\ref{sec:ontology} derives the constraints that any internal-rotation ontology of spin must satisfy,
and reviews the known constructions that meet them.
Section~\ref{sec:causality} shows how the rotating-wave formulation dissolves the historical
superluminal-velocity objection while distinguishing internal rotation from the separate causality problem
raised by any literal spatial flow.
Section~\ref{sec:cosmochrony} records the consequences for the Cosmochrony programme.
Section~\ref{sec:higgsboundary} reports an independent downstream reconnaissance: an exact
representation audit showing that fermionic chiral bilinears carry the Higgs doublet's quantum
numbers, without deriving a composite condensate or its uniqueness.
Section~\ref{sec:audit} reports the outcome of a soldering audit of those consequences: an unconditional
Veronese-type obstruction lemma and a corpus-conditional indeterminacy statement.
Section~\ref{sec:conclusion} states the resulting structural hierarchy.

\section{The photon: rotating fields and helicity}
\label{sec:photon}

In a plane electromagnetic wave in vacuum the electric and magnetic fields are perpendicular in space but
oscillate in phase: $\mathbf{B} = \hat{\mathbf{k}} \times \mathbf{E}/c$, and both vanish simultaneously.
The rotation relevant to spin does not involve a phase shift between $\mathbf{E}$ and $\mathbf{B}$.
It involves two transverse components of the electric field itself,
\begin{equation}
\mathbf{E}(z,t) = E_x \hat{\mathbf{x}} \cos\alpha + E_y \hat{\mathbf{y}} \cos(\alpha + \delta),
\qquad \alpha = kz - \omega t + \alpha_0 .
\label{eq:transverse}
\end{equation}
For $\delta = \pm\pi/2$ and $E_x = E_y$ the field vector rotates about the propagation axis, and the sense
of rotation corresponds to the helicity $\lambda = \pm 1$ of the photons, i.e.\ to the projection
$S_\parallel = \pm\hbar$ of their spin on the propagation direction.

Two facts support the rotational intuition here.
First, the angular momentum carried by circularly polarized light is mechanically real: a suspended
birefringent plate experiences a measurable torque when it reverses the polarization of a light
beam~\cite{Beth1936}.
Second, the phase relation $\delta$ between the transverse components does determine the polarization state
and hence the mean helicity.

Two facts, however, bound the intuition.
Linear polarization ($\delta = 0$ or $\pi$) is not a zero-helicity photon state; it is a coherent
superposition of the two helicities, and a helicity measurement yields $\pm\hbar$, never $0$.
The absence of a longitudinal (helicity-zero) mode for the free photon follows from masslessness and gauge
invariance, through the finite-helicity representations of the massless little group $ISO(2)$, in which the
translational part acts trivially~\cite{Wigner1939}; Wigner's classification also contains continuous-spin
representations, and a massive spin-$1$ boson does possess three states.
Moreover, a one-photon state of definite helicity has a vanishing mean electric field,
$\langle \mathbf{E} \rangle = 0$, whereas a circularly polarized coherent state has a rotating mean field.
The rotating classical field is therefore the coherent-state manifestation of the helicity structure, not a
literal microscopic mechanism of photon spin.

\section{Rotating waves rather than rotating objects}
\label{sec:rotatingwaves}

The founding hypothesis can now be stated precisely: the rotating entity is not a particle-like object but
a wave mode, the wave function being a combination of rotating (chiral) and non-rotating components.

For the photon this language is physically literal.
The state decomposes on the two circular polarization modes,
$\Psi = \psi_+ \ket{+} + \psi_- \ket{-}$, which are the two helicity sectors, and the Maxwell fields can be
packaged into the Riemann--Silberstein combinations
\begin{equation}
\mathbf{F}_\pm = \mathbf{E} \pm i c \mathbf{B},
\label{eq:riemannsilberstein}
\end{equation}
which separate the sectors of opposite helicity and support a genuine photon wave-function
formalism~\cite{BialynickiBirula1996}.
For real Maxwell fields $\mathbf{F}_- = \overline{\mathbf{F}_+}$, so the two combinations are not
independent real fields; the clean helicity separation is obtained after complexification and restriction
to the positive-frequency components.
Speaking of right- and left-rotating waves is then physically pertinent.
A non-rotating wave, such as a linear polarization, is not a spin-zero state but a coherent superposition of
the two counter-rotating modes; its mean angular momentum may vanish, yet a helicity measurement always
yields $\pm\hbar$.
The distinction $\langle S \rangle = 0$ versus $j = 0$ is essential: the former can result from the mutual
cancellation of two opposite modes, the latter designates a wave that transforms as a genuine scalar under
rotations.

A counterexample fixes the decisive criterion.
A complex scalar wave $\psi = A e^{-i\omega t}$ can be viewed as two real components in quadrature,
$\psi = A\cos\omega t - iA\sin\omega t$, and thus ``rotates'' in its complex plane; yet its quanta can be of
spin zero.
The rotation of two components is therefore not sufficient: what matters is how the components transform
under rotations of physical space,
\begin{equation}
\Psi(\mathbf{x}) \longmapsto D^{(j)}(R)\, \Psi(R^{-1}\mathbf{x}).
\label{eq:transformationlaw}
\end{equation}
If $D(R) = 1$ the wave is scalar (spin $0$); a vectorial $D(R)$ gives spin $1$; a spinorial $D(R)$ changing
sign after $2\pi$ gives spin $\tfrac12$; higher representations give $j = \tfrac32, 2, \ldots$
The internal rotation counts only insofar as its plane is geometrically coupled to the directions of
physical space; a rotation in a purely internal plane, as in the scalar counterexample, carries no spin.

The rotating-wave hypothesis can thus be stated rigorously: a one-particle state is represented by a
carrier-valued wave amplitude; its spin is fixed by the carrier's transformation law; its spin state
indicates which combination of waves rotating in the various senses is realized when such a decomposition
exists.
For a massive fermion, however, chirality, helicity, and the weight $m$ are distinct notions, and the
counter-rotating decomposition is literal only in special cases such as photon polarization; in general the
components form a carrier space whose transformation law, not its visual rotation, defines the spin.
This is compatible with current physics, and it is native to the ground where the deeper constructions of
Section~\ref{sec:ontology} operate, since solitons, skyrmions, and Finkelstein--Rubinstein quantization
concern extended field configurations rather than point objects.
What remains to be derived is why the components of the wave carry precisely a given representation
$D^{(j)}$: the rotation of the wave can represent the spin, but the topology and transformation law of the
wave must still explain why it is scalar, vectorial, or spinorial.

Finally, this internal rotation must be distinguished from a vortex wave,
$\psi(r, \varphi, z) \propto e^{i\ell\varphi}$, whose phase genuinely winds around a spatial axis and
produces an orbital angular momentum $L_z = \ell\hbar$.
A wave can carry both, $J_z = L_z + S_z$, with a spatial orbital winding and an internal polarization
rotation.
For the electromagnetic field the split of $J$ into orbital and spin parts is not gauge invariant in full
generality; it becomes operational in the paraxial regime, where spin and orbital angular momenta of
Laguerre--Gaussian beams are separately measurable~\cite{Allen1992}.
The rotating-wave intuition targets the spin part.

\section{Weights, windings, and representations}
\label{sec:representations}

The temptation is to associate the discrete spin values $j = 0, \tfrac12, 1, \tfrac32, \ldots$ with
particular phase angles such as $\delta = 0$, $\pi/2$, $\pi$ in Eq.~\eqref{eq:transverse}.
This fails, because $\delta$ varies continuously and parametrizes states \emph{within} a given state space,
whereas $j$ labels the space itself through its transformation law under all rotations.

A physical rotation acts on a spin-$\tfrac12$ state by
\begin{equation}
U(\mathbf{n}, \alpha) = \exp\!\left(-\frac{i\alpha}{2}\, \mathbf{n} \cdot \boldsymbol{\sigma}\right),
\label{eq:su2rotation}
\end{equation}
so that $U(2\pi) = -I$ and only $U(4\pi) = +I$.
The topology of the rotation group, $\pi_1(SO(3)) = \mathbb{Z}_2$, explains this dichotomy: it allows
exactly two behaviours under a $2\pi$ rotation, $U(2\pi) = \pm I$, hence the integer and half-integer
families.
But the precise values of $j$ are not elements of $\pi_1$; they arise from the theory of irreducible
representations of $SU(2)$, where $2j$ is an integer highest weight, the degree of the representation
$\operatorname{Sym}^{2j}(\mathbb{C}^2)$.

A genuine winding number does survive after restriction to rotations about a fixed axis:
\begin{equation}
U_z(\alpha)\, \ket{j, m} = e^{-i m \alpha}\, \ket{j, m},
\label{eq:weight}
\end{equation}
where $2m \in \mathbb{Z}$ is the winding degree of this $U(1)$ character over a full cycle of the covering
group ($4\pi$).
If $m$ is half-integer, $e^{-i m (2\pi)} = -1$ and $e^{-i m (4\pi)} = +1$.
The winding intuition is thus rigorous but axial and local: it describes the weights $m$, and $j$ appears
only when these weights are organized into a complete $SU(2)$ representation.
It is also through the finite-helicity representations of the little group, together with the global choice
of the relevant covering for half-integer values, that the helicity of massless particles is
quantized~\cite{Wigner1939}.

For a spin-$\tfrac12$ state written as
\begin{equation}
\ket{\psi} = \cos\frac{\theta}{2}\, \ket{\uparrow} + e^{i\phi} \sin\frac{\theta}{2}\, \ket{\downarrow},
\label{eq:bloch}
\end{equation}
the continuous relative phase $\phi$ determines the orientation of the state on the Bloch sphere, not the
intrinsic value of the spin: a measurement along a fixed axis yields only $\pm\hbar/2$.

\section{Interferometric observability of the spinor sign}
\label{sec:interferometry}

The sign reversal $U(2\pi) = -I$ is a global phase if the whole state is rotated, and as such is
unobservable.
It becomes observable when one branch of an interferometer is rotated relative to the other.
Neutron interferometry realizes this: after a relative $2\pi$ spin rotation the two amplitudes interfere
with opposite sign, and the initial interference pattern is restored only after
$4\pi$~\cite{Rauch1975, Werner1975}.

What the experiment detects is not a hidden orientation of the neutron, but the spinorial character of the
action of rotations on its quantum amplitude.
The physically real structure is the transformation class of the state under rotations, since a state
``rotated by $2\pi$'' can be made to interfere with itself and the difference observed.
The rotation phase itself, by contrast, remains unobservable in isolation.

\section{Univalence as a rule on the algebra of observables}
\label{sec:univalence}

Since $U(2\pi) = +1$ on integer-spin states and $-1$ on half-integer-spin states, a coherent superposition
of the two would change its relative phase under a $2\pi$ rotation, that is, under no physical operation at
all.
The univalence superselection rule~\cite{Wick1952} draws the consequence.

The rule must be stated carefully.
It does not forbid writing the vector sum $\ket{\psi_{\mathrm{B}}} + \ket{\psi_{\mathrm{F}}}$ of a bosonic
and a fermionic state; such a vector exists mathematically.
It states that no admissible physical observable measures its relative phase.
Indeed, a $2\pi$ rotation maps
\begin{equation}
\ket{\psi_{\mathrm{B}}} + e^{i\phi} \ket{\psi_{\mathrm{F}}}
\;\longmapsto\;
\ket{\psi_{\mathrm{B}}} - e^{i\phi} \ket{\psi_{\mathrm{F}}},
\label{eq:univalence}
\end{equation}
and these two descriptions must be physically indistinguishable.
The phase $\phi$ therefore enters no observable, and the would-be superposition is operationally equivalent
to the corresponding mixture.
In relativistic quantum field theory this becomes essentially the superselection of fermion parity.

Two kinds of phase must accordingly be distinguished.
Within a single sector, as in neutron interferometry where both branches belong to the same half-integer
sector, relative phases are measurable.
Between the integer and half-integer sectors, no physical observable can serve as a coherent bridge, and the
relative phase is suppressed by superselection.
Univalence is thus a rule about the algebra of observables, not a mathematical impossibility of writing
vector sums.

\section{Constraints on internal-rotation ontologies}
\label{sec:ontology}

The preceding results define a precise specification for any interpretation of spin as an underlying
rotation.
Such a construction should not merely exhibit something that rotates.
It should explain why $R(2\pi) \neq I$ while $R(4\pi) = I$ on the underlying object, render the first effect
observable only by relative comparison, and reproduce the measurement probabilities, relativistic
covariance, and the integer versus half-integer dichotomy.

A simple rotating arrow fails this test: an arrow returns to itself after $2\pi$.
The minimal classical configuration with the required topology is exhibited by Dirac's belt trick: a frame
attached to its environment, whose configuration retains the relation to its surroundings or to the history
of its rotation.
A $2\pi$ rotation of such a frame cannot be undone with the attachments held fixed, whereas a $4\pi$
rotation can.
This is the ``orientation-entanglement'' of Misner, Thorne, and Wheeler~\cite{Misner1973}.

Orientation-entanglement is not, however, the only construction available.
Extended objects can acquire fermionic behaviour upon quantization: solitons and skyrmions~\cite{Skyrme1961}
under the Finkelstein--Rubinstein quantization conditions~\cite{Finkelstein1968}, and gravitational geons in
the same framework.
Their common feature confirms the structural lesson: the spinorial sign never comes from a local arrow, but
from the global topology of the configuration space.

Closest to the initial intuition, Hestenes represents the Dirac spinor as a rotor of the Clifford algebra,
encoding a local frame, a phase, and a rotation plane, in connection with the zitterbewegung of Dirac
theory~\cite{Hestenes1990}.
This is a substantial and illuminating geometric reformulation of the Dirac formalism, but it does not by
itself derive spin $\tfrac12$ from a more fundamental substructure.

The minimal ontological hypothesis compatible with all constraints is therefore the following: if spin
originates in an underlying rotation, the rotating entity must be a relational, extended, or topologically
attached wave configuration, never a point particle endowed with a mere internal axis.
The rotating-wave formulation of Section~\ref{sec:rotatingwaves} meets the negative part of this
requirement, since it discards the rigid point-like object; but it does not by itself derive
$U(2\pi) = -I$, the spinorial topology, or the quantum measurement probabilities, which remain precisely
the results a deeper construction would have to produce.

\section{Causality and the superluminal-velocity objection}
\label{sec:causality}

Any rotational model of the electron faces a famous objection: if a sphere of classical electron radius
carries an angular momentum $\hbar/2$, its equatorial velocity must exceed the speed of light by orders of
magnitude.
The difficulty was already noted by Uhlenbeck and Goudsmit themselves~\cite{Uhlenbeck1925}, formulated by
Lorentz, and later pressed by Pauli against any material interpretation of spin; the detailed history is
given by Giulini~\cite{Giulini2008} (see also~\cite{Tomonaga1997}).
The objection targets a specific model: matter distributed at a spatial radius $r$, rotating mechanically at
angular frequency $\omega$, with local velocity $v = \omega r$.
Three quantities must therefore be kept distinct: the angular frequency of an internal evolution, the phase
velocity of a wave pattern, and the tangential velocity of a material distribution.
Only the third is directly constrained by relativity.

The standard formalism, moreover, assigns the electron no intrinsic mechanical rotation frequency at all.
A stationary state of definite spin evolves as $\ket{\psi(t)} = e^{-iEt/\hbar} \ket{\mathbf{p}, s}$: the
global phase rotates at frequency $E/\hbar$ while $\langle \mathbf{S} \rangle$ remains constant.
Spin precession occurs only in an external magnetic field, at the Larmor frequency
$\boldsymbol{\Omega}_L \simeq -(gq/2m)\, \mathbf{B}$, which describes the evolution of the quantum
orientation, not the rotation of a material surface.

For the rotating-wave hypothesis of Section~\ref{sec:rotatingwaves} the objection dissolves.
A rotation in the internal space of the components, $\Psi \mapsto U(t)\, \Psi$, has for ``radius'' a field
amplitude, not a distance; the product $\omega r$ is then not a spatial velocity.
Circular polarization illustrates this: at a fixed point the direction of $\mathbf{E}$ rotates at the
optical frequency, yet no material surface traverses a circle, and in an ideal plane wave the energy flux
remains directed along the propagation.

Causality nevertheless remains binding wherever a literal circulation is asserted.
If a density of charge or energy really circulates around a spatial axis, its local velocity must respect
relativity; for the electromagnetic field the energy flux obeys $|\mathbf{S}_P| \le c\, u$, with
$\mathbf{S}_P$ the Poynting vector and $u$ the energy density.
A phase pattern may travel faster than $c$ without transporting information or energy at that speed, but a
real flow may not.

The wave formulation admits a field-theoretic reading of this constraint.
Following Belinfante~\cite{Belinfante1939}, Ohanian showed that the Dirac angular momentum can be expressed
in terms of a circulating energy flow in the wave field, extended over the Compton scale
$\lambda_C = \hbar / mc$~\cite{Ohanian1986}.
This removes the need for a surface velocity $\omega r$ at any assigned radius: the classical paradox arose
from confining the rotation of a rigid body to the classical electron radius
$r_e = \alpha \lambda_C \ll \lambda_C$.
The result must not be over-read, however.
It establishes an equivalent field-theoretic bookkeeping of the angular momentum as a flow, not literal
circular trajectories with a unique local velocity, and the electromagnetic bound
$|\mathbf{S}_P| \le c\, u$ does not transpose automatically to the Dirac energy--momentum tensor.
A literal hydrodynamic interpretation of this circulation would therefore require a separate causal
analysis.

The verdict is thus twofold.
An internal rotation of the wave components raises no superluminal problem at all.
A spatial-circulation reading is available at the level of the field's energy flow on the Compton scale,
but its literal hydrodynamic interpretation is not established by the formalism, and the rotating-wave
hypothesis must refrain from reintroducing, surreptitiously, a mechanical radius of the electron.

\section{Implications for the Cosmochrony programme}
\label{sec:cosmochrony}

The preceding structure sorts the current state of the Cosmochrony construction~\cite{Cosmochrony} into
established facts, one exact consequence, one needed correction, one diagnostic, and one open question.

Established: the admissibility construction selects a two-dimensional complex module
$V_\rho \simeq \mathbb{C}^2$,
with $\operatorname{Sym}^2(V_\rho) = V_1$ and $\wedge^2(V_\rho) = V_0$, reproducing the Clebsch--Gordan
content $\tfrac12 \otimes \tfrac12 = 0 \oplus 1$.

Exact algebraic consequence: the central element $-I$ acts on $V_\rho^{\otimes n}$ as $(-1)^n$.
Conditional on the Lorentz-spin identification of $V_\rho$, this tensor grading becomes the physical
univalence grading of Section~\ref{sec:univalence}; absent that identification, it remains a grading of the
admissible $SU(2)$ only.

Correction: in the Q7 construction~\cite{Beau2026q7}, the state $e_0$ must be identified as the $m = 0$
weight of the $j = 1$ module, not as a $j = 0$ state; the genuine $j = 0$ state is the antisymmetric
singlet in $\wedge^2(V_\rho)$.

Diagnostic: Lorentzian spin and the weak doublet must act on distinct tensor factors; a single copy of
$V_\rho$ cannot carry both actions without conflating spacetime and internal symmetries.

Open questions: two distinct problems must be separated.
The representational problem is to construct the full action
$\rho_{\mathrm{rot}} \colon \operatorname{Spin}(3) \to U(\mathcal{H}_\Pi)$ with
$\rho_{\mathrm{rot}}(-I) = (-1)^n I$ on the tensor sector of degree $n$; a one-parameter evolution cannot
supply the three non-commuting generators of $\operatorname{Spin}(3)$, so the cascade alone cannot solve
this problem.
The dynamical problem, stronger, is to determine whether the cascade canonically traces the orbit of an
axial one-parameter subgroup,
$U_\Pi(\tau) = \rho_{\mathrm{rot}}\!\left(\exp[-i \alpha(\tau) J_{\mathbf{n}}]\right)$, with both the axis
$\mathbf{n}$ and the relation $\alpha(\tau)$ derived rather than presupposed.
Solving the first would be a necessary step toward a rotational derivation of spin; solving the second
would tie that representation to the cascade dynamics; neither is yet a derivation of the full quantum
formalism.
These two problems have since been subjected to a systematic audit against the deposited corpus; its
formal outcomes are recorded in Section~\ref{sec:audit}.
A separate, independent downstream reconnaissance, reported next in Section~\ref{sec:higgsboundary},
asks whether the same chiral structure typed by Eq.~\eqref{eq:transformationlaw} has a further
physical role: whether it is what lets a fermionic wave couple to the Higgs condensate.

\section{Downstream boundary: chiral transversality and Higgs-typed composite channels}
\label{sec:higgsboundary}

The transformation-law criterion of Eq.~\eqref{eq:transformationlaw} already fixes what carries spin.
A natural downstream question is whether the same chiral structure is what lets a fermionic wave
couple to the Higgs condensate, as opposed to some other, non-interacting spin assignment.
Stated loosely, this conflates several notions this note has kept apart (spin, helicity, chirality).
Stated precisely, it asks whether the two Lorentz-chiral components $S_L, S_R$ of a spin-$\tfrac12$
carrier admit a transverse operator, supplied by the Higgs condensate, converting one into the other,
with mass measuring the norm of that inter-chiral incidence.
This section reports an exact reconnaissance of that question, entirely downstream of, and
independent from, the soldering audit of Section~\ref{sec:audit}: it neither uses nor repairs the
missing rotational soldering $\rho_{\mathrm{sp}}$, and its results hold whether or not that soldering
is ever supplied.

\subsection{The chiral Hamiltonian and Bloch-vector precession}
\label{sec:chiralhamiltonian}

For fixed helicity $\eta = \pm 1$, in the chiral basis $(\psi_L, \psi_R)$, the standard electroweak
Yukawa coupling~\cite{Weinberg1967, Higgs1964} gives the two-level Hamiltonian
\begin{equation}
H_\eta = \mathbf{B}_\eta \cdot \boldsymbol{\tau}, \qquad
\mathbf{B}_\eta = \bigl(\operatorname{Re} m,\, -\operatorname{Im} m,\, -\eta |\mathbf{p}|\bigr),
\label{eq:chiralham}
\end{equation}
with $\boldsymbol{\tau}$ the Pauli matrices acting on the chiral index and
$m = (y/\sqrt2)(v + h)$ the position-dependent Yukawa mass term.
Verified by exact symbolic computation (script cited in Section~\ref{sec:reproducibility}):
\begin{equation}
H_\eta^2 = \bigl(|\mathbf{p}|^2 + |m|^2\bigr) I, \qquad
\text{eigenvalues } \pm\sqrt{|\mathbf{p}|^2 + |m|^2},
\label{eq:chiralspectrum}
\end{equation}
and the chiral pseudo-spin $\mathbf{s} = \langle \boldsymbol{\tau} \rangle$ precesses as
\begin{equation}
\dot{\mathbf{s}} = 2\, \mathbf{B}_\eta \times \mathbf{s},
\label{eq:chiralprecession}
\end{equation}
an exact operator identity, $i[H_\eta, \tau_i] = 2 (\mathbf{B}_\eta \times \boldsymbol{\tau})_i$ for
each $i$, verified symbolically rather than in a particular representation.
At $m = 0$ the chiralities decouple exactly.
Under relative chiral rephasing $\psi_L \mapsto e^{i\alpha} \psi_L$, $\psi_R \mapsto e^{i\beta} \psi_R$,
$m \mapsto m\, e^{i(\beta - \alpha)}$, so $H_\eta$ transforms into itself with $m$ rotated: $\arg m$ is
the basis-dependent orientation of the transverse chiral component, removable for a single Dirac
sector by this rephasing.
Only a phase invariant under every such single-sector rephasing, necessarily built from a relative
comparison across several Yukawa matrices such as a CKM-type phase, can be physical.
This precession is the exact content behind ``the Higgs condensate as a transverse chiral
direction'': a statement about the chiral Bloch vector $\mathbf{s}$, never about a spatial rotation
of the wave, and it must not be conflated with the transformation law of
Eq.~\eqref{eq:transformationlaw} that fixes spin itself.

\subsection{Higgs-typed composite channels: exact representation content}
\label{sec:compositegate}

The chiral picture above reproduces known electroweak structure and says nothing new by itself.
The open question is representational: does the Dirac-conjugated bilinear $\overline{\psi_R} \psi_L$
carry the Higgs doublet's quantum numbers, so that a dynamical condensate of that bilinear, in the
spirit of the dynamical-symmetry-breaking mechanism of~\cite{NambuJonaLasinio1961} and of the
top-quark-condensate construction of~\cite{BardeenHillLindner1990}, could in principle realise an
effective composite Higgs doublet?
This subsection answers only the representation-theoretic half of that question, for each tree-level
Yukawa sector ($e$, $d$, $u$), exactly, under
$\operatorname{Spin}(3,1) \times SU(3)_c \times SU(2)_L \times U(1)_Y$:
\begin{itemize}
\item \textbf{Lorentz scalar:} writing the infinitesimal $SL(2,\mathbb{C})$ action on left/right
  two-component spinors with a shared rotation generator and opposite-sign boost
  generators~\cite{DreinerHaberMartin2010}, the operator identity
  $\mathrm{gen}_R^\dagger + \mathrm{gen}_L = 0$ holds exactly for arbitrary rotation and boost
  parameters: $\psi_R^\dagger \psi_L$ is Lorentz-invariant.
\item \textbf{Color:} the quark channels $\overline{d}_R Q_L$ and $\overline{u}_R Q_L$ are exact
  $SU(3)_c$ singlets, verified as an exact invariance identity under all eight Gell-Mann generators.
\item \textbf{Weak isospin and hypercharge:} the $SU(2)_L$ content is the trivial doublet
  $1 \otimes 2 = 2$; exact hypercharge arithmetic gives
  \begin{equation}
  \begin{array}{c|c|c}
  \text{composite channel} & Y & \text{Higgs type} \\ \hline
  \overline{e}_R L_L & +\tfrac12 & H \\
  \overline{d}_R Q_L & +\tfrac12 & H \\
  \overline{u}_R Q_L & -\tfrac12 & \widetilde{H}
  \end{array}
  \label{eq:higgstable}
  \end{equation}
  with $\widetilde{H} = i \sigma_2 H^*$ the conjugate doublet.
\item \textbf{Electroweak breaking:} each channel has an exact electrically neutral slot
  ($Q = T_3 + Y = 0$) annihilated by the unbroken $U(1)_{\mathrm{em}}$ generator.
\end{itemize}
All identities above are verified by exact symbolic computation with no sampling.

\subsection{Stratified verdict}
\label{sec:higgsverdict}

Every channel is separately Higgs-typed; no computation here supplies a canonical map identifying the
three channels, built respectively from the lepton doublet and from the two quark chiralities, as
conjugate realisations of one composite order.
The audited gauge group $\operatorname{Spin}(3,1) \times SU(3)_c \times SU(2)_L \times U(1)_Y$
contains no generator relating the lepton sector to either quark sector: the three channels are
independently typed, not images of one another under that group.
\begin{equation}
\boxed{
\begin{aligned}
\text{spinorial chirality} &\;\Longrightarrow\; \text{Higgs-typed composite channels are allowed},\\
\text{spinorial rotation} &\;\not\Longrightarrow\; \text{a unique Higgs condensate is generated}.
\end{aligned}}
\label{eq:higgsverdict}
\end{equation}
Three data remain open, and must not be collapsed into one.
(i) The composite-channel construction of Section~\ref{sec:compositegate} requires each fermion to
simultaneously carry Lorentz-chiral, electroweak, and (for quarks) color structure: exactly the
ambient enlargement beyond a single copy of $V_\rho$ already flagged as required, and not yet
supplied, by the diagnostic of Section~\ref{sec:cosmochrony}.
(ii) A canonical mechanism relating or selecting the three channels as a single common order is
absent from the audited gauge content.
(iii) The condensation dynamics itself, and the selection of the electroweak scale, remain
undetermined.
Conditional on (i), this audit leaves (iii) as the next obstruction to producing at least one
composite Higgs channel; producing one universal Higgs order shared by all three Yukawa sectors
additionally requires (ii).
Opening any of these three is a new front, materially larger than this reconnaissance, and is not
undertaken here.

\subsection{Reproducibility}
\label{sec:reproducibility}

All identities of this section are verified by an exact symbolic script,
\texttt{code/\allowbreak front\_\allowbreak spinorial\_\allowbreak higgs\_\allowbreak
transversality.py} (pinned \texttt{sympy==1.14.0}, no sampling), documented and run from a fresh
environment as described in the README.

\section{The soldering audit: a Veronese obstruction and an indeterminacy statement}
\label{sec:audit}

The open problems of Section~\ref{sec:cosmochrony} were audited against the deposited corpus under a
pre-registered kill-switch: if the identification of the admissible $SU(2)$ with the double cover of the
emergent rotations requires a free choice of basis, orientation, compact real form, or a non-derived
soldering morphism, then the programme does not yet derive physical spin.
This section records the two formal outcomes: an unconditional obstruction lemma and a corpus-conditional
indeterminacy statement.

\subsection{Classification of the real forms of the carrier square}
\label{sec:realforms}

Let $V_\rho \simeq \mathbb{C}^2$ carry the admissible $SU(2)$ action~\cite{Beau2026a33}, let $J$ be its
quaternionic structure (unique up to a phase), and set $J_\phi = e^{i\phi} J$.

\begin{proposition}[Real forms of $\operatorname{Sym}^2(V_\rho)$]
\label{prop:realforms}
(i) Each induced map $J_\phi^{(2)}$ on $\operatorname{Sym}^2(V_\rho)$ is an antilinear involution, with
fixed real form $W_\phi = e^{i\phi} W_0$, where
$W_0 = \operatorname{span}_{\mathbb{R}}\{x_1 + x_3,\; i(x_1 - x_3),\; i x_2\}$ in the basis
$\{x_1 = e_1 e_1,\, x_2 = e_1 e_2,\, x_3 = e_2 e_2\}$.
(ii) The $W_\phi$ form a circle of distinct real three-dimensional forms, permuted by the center of
$U(2)$ and not by $SU(2)$; each is $SU(2)$-invariant.
(iii) The induced real action on each $W_\phi$ is $SO(3)$, with kernel exactly $\{\pm I\}$.
(iv) Neither the Hermitian metric of the unitary structure nor the invariant bilinear form $B$, defined
only up to $\mathbb{C}^*$, distinguishes the $W_\phi$: the rescaled form $B_\phi = e^{-2i\phi} B$
restores the same positive Gram matrix $\operatorname{diag}(2, 2, \tfrac12)$ on every $W_\phi$.
\end{proposition}

\begin{proof}
All statements reduce to polynomial identities on the Cayley--Klein chart, verified by direct symbolic
computation: $U E = E \overline{U}$ for every $U \in SU(2)$ with $E$ the matrix of $J$, hence
$\operatorname{Sym}^2(U E) = \operatorname{Sym}^2(E \overline{U})$ (invariance of every $W_\phi$,
without using the unit-norm constraint); $S^T B S = (|a|^2 + |b|^2)^2 B$ identically for
$S = \operatorname{Sym}^2(U)$ (invariance of $B$ on $SU(2)$); and $\operatorname{Sym}^2(U) = I$ forces
$b = 0$ and $a = \pm 1$, so the kernel equals $\{\pm I\}$ rather than merely containing it.
For surjectivity in (iii): since $B_\phi$ is positive definite on $W_\phi$, the image of the induced
action lies in $O(3)$; connectedness of $SU(2)$ places it in $SO(3)$; and the kernel being discrete, the
image is three-dimensional, hence all of $SO(3)$.
\end{proof}

The residual central $U(1)/\{\pm 1\}$ ambiguity is therefore intrinsic to the data
$(V_\rho, SU(2), h, B)$: a real-form selector must be a corpus-derived, phase-sensitive datum.

\subsection{A Veronese obstruction to the first internal selector}
\label{sec:veronese}

The natural internal candidate for such a selector is the family of per-pair outer products
$M_j = w_j w_j^{\mathsf{T}}$ that spans $\mathcal{H}_{\mathrm{eff}} \simeq
\operatorname{Sym}^2(V_\rho)$~\cite{Beau2026a32}: if their real span coincided with a single $W_\phi$,
the phase would be fixed by projection data.
This candidate is excluded exactly.

\begin{lemma}[Veronese-null versus compact real form]
\label{lem:veronese}
For $w = (u, v) \in \mathbb{C}^2$, the pure tensor $M = w w^{\mathsf{T}}$ has coordinates
$(u^2,\, 2uv,\, v^2)$ and satisfies $B(M, M) = 2 u^2 v^2 - \tfrac12 (2uv)^2 = 0$: every nonzero pure
tensor lies on the complex null cone of $B$ (the Veronese cone).
Since every $W_\phi$ is positive definite for $B_\phi$, it meets the null cone only in $0$.
\end{lemma}

\begin{corollary}
\label{cor:noselector}
$\operatorname{span}_{\mathbb{R}}\{M_j\} \neq W_\phi$ for every $\phi$, even when that real span has
dimension exactly three: the outer products of~\cite{Beau2026a32} cannot select the real form.
\end{corollary}

\begin{proof}
Each $M_j$ lies in $\operatorname{span}_{\mathbb{R}}\{M_j\}$ and on the null cone; a nonzero null vector
cannot lie in a subspace on which the (rescaled) form is definite.
\end{proof}

\subsection{Corpus inventory and the indeterminacy statement}
\label{sec:indeterminacy}

The audit traced the remaining soldering inputs through the deposited corpus.
First, the normalisation: the spatial coefficient $A_Z = 2$ is set equal to the $j = 1$ Casimir
eigenvalue in~\cite{BeauQ8} and inherited by~\cite{Beau2026q10} and~\cite{Beau2026q11}, whereas the
isotropic \emph{shape} of the effective metric is derived independently (Born--Infeld parity, the
universality theorem [U] of~\cite{Beau2026u1}, and the absence of cross terms proved
in~\cite{Beau2026q7}).
The identity ``effective metric $=$ Casimir'' is therefore currently definitional rather than a match;
a Casimir-independent normalisation (for instance the announced sub-principal-symbol route
of~\cite{BeauQ8}) remains to be executed.
Second, the target action: the deposited corpus defines no rotation action on the three effective
directions independent of the admissible module — the concrete symmetries are the $U(1)$ generated by
$J_3$ and the metaplectic horizontal rotation, and the rotational structure of the effective space is
transported through the bridge of~\cite{Beau2026q7}, conditionally on that bridge.
Third, orientation is \emph{not} blocking: the chiral orientation result
[H-orient]~\cite{Beau2026cho, Beau2026aog}, established within the present spin stratum by purely
representation-theoretic and arithmetic means, uses no soldering morphism and can therefore orient a
future one without circularity.

The outcome is a conditional statement, relative to the inventory just listed.

\begin{proposition}[Indeterminacy of the spinorial soldering]
\label{prop:indeterminacy}
The currently derived data — the admissible module $V_\rho$ with its unitary structure and invariant
form, the central grading $(-1)^n$ on $V_\rho^{\otimes n}$, the isotropic shape of the effective metric,
and the horizontal symmetries — determine the complex carrier and its univalence, but are invariant
under the residual family of solderings parametrised by the circle $\{W_\phi\}$ and admit no
independently defined intertwining target.
They therefore determine neither the spatial realisation of the carrier, nor a physical rotation action,
nor its scale.
\end{proposition}

The missing soldering thus comprises three independently missing data: a phase-sensitive real-form
selector (any real-span candidate must supply derived vectors of nonzero $B$-norm, by
Lemma~\ref{lem:veronese}), an independent rotation action $\rho_{\mathrm{sp}}$ on the effective space,
and a Casimir-independent normalisation.
Whether these three data would jointly suffice for a complete soldering is a separate sufficiency
question, not settled here.
Their priority order is structural: the absence of $\rho_{\mathrm{sp}}$ is irrecoverable by the other
two.
Lemma~\ref{lem:veronese} is unconditional mathematics; Proposition~\ref{prop:indeterminacy} is
conditional on the listed corpus inventory: it is falsified if any of these data is shown to be already
derivable from that inventory, and it must be reassessed whenever a new datum is derived.

\section{Conclusion}
\label{sec:conclusion}

The founding hypothesis, spin as the rotation of the wave modes composing the wave function, survives in a
corrected and stratified form:
\begin{equation}
\boxed{
\begin{aligned}
\text{relative phase} &: \ \text{measurable coordinate within a sector},\\
m &: \ \text{winding weight of the rotation about an axis},\\
j &: \ \text{global representation of } SU(2),\\
(-1)^{2j} &: \ \text{univalence separating the physical sectors}.
\end{aligned}}
\label{eq:hierarchy}
\end{equation}
Spin is not the instantaneous angle of an internal rotation; it characterizes the global way the wave mode
transforms under rotations, as expressed by Eq.~\eqref{eq:transformationlaw}.
A one-particle state is represented by a carrier-valued wave amplitude, and its spin is fixed by the
carrier's transformation law; only in special cases, notably photon polarization, can its components be
interpreted literally as counter-rotating modes.
For the photon this structure manifests itself classically as the chirality of polarization, measurable as
mechanical angular momentum~\cite{Beth1936}.
For fermions it requires a spinorial structure whose relative phases have measurable interferometric
consequences~\cite{Rauch1975, Werner1975}, and whose sign structure partitions the Hilbert space into
superselection sectors~\cite{Wick1952}.
The hypothesis escapes the superluminal-velocity objection, which constrains rigid bodies, not internal
rotations of wave components~\cite{Ohanian1986}.
The hierarchy of Eq.~\eqref{eq:hierarchy} is thus an experimentally anchored structure fixing the
specification that any deeper rotational ontology of spin must meet.
The soldering audit of Section~\ref{sec:audit} turns that specification into a falsifiable programme
statement: the deposited construction owns the abstract spinorial carrier, its exact double cover, and
its univalence grading, and provably does not yet own physical spin — the rotational soldering is
missing and is typed as three data, with the independent rotation action $\rho_{\mathrm{sp}}$ as the
structurally blocking one.
The founding hypothesis is thereby neither confirmed nor vague: each failed candidate now produces a
lemma, as Lemma~\ref{lem:veronese} illustrates.
A separate downstream reconnaissance (Section~\ref{sec:higgsboundary}) shows that the same chiral
structure typed by Eq.~\eqref{eq:transformationlaw} is compatible with, but does not derive, a further
physical role: fermionic chiral bilinears carry exactly the Higgs doublet's quantum numbers in every
audited Yukawa sector, while neither the physical rotation of the wave, the gauge group, nor this
representation typing determines whether, which, or at what scale such a bilinear condenses.

\subsection*{Interpretive outlook}

The following is an interpretive reading, not a consequence of the preceding propositions.
If the missing soldering data were derived and shown to be jointly sufficient, the resulting
construction could support an emergent reading of spatial rotation: rotations would then act as
effective symmetries of projected geometry, and the spinor sign could be interpreted as a topological
memory of projection rather than as an independently imposed representation-theoretic input.
Conversely, if a stronger no-go result established that the residual indeterminacy is irreducible, this
would be consistent with --- but would not by itself prove --- a spinorial carrier prior to emergent
space, with univalence inherited by the latter.
The present audit does not decide between these readings.
Its result is to locate the exact boundary between what the present construction derives and what it
still postulates.
It therefore neither confirms nor refutes the rotating-wave hypothesis; it turns its realisation within
Cosmochrony into an explicit derivation problem with testable success and failure criteria.
The downstream boundary of Section~\ref{sec:higgsboundary} is independent of the soldering question
addressed here: it neither uses nor repairs the missing rotational soldering $\rho_{\mathrm{sp}}$, and
holds under either of the two readings above.

\clearpage
\bibliographystyle{plain}
\bibliography{cosmochrony-bibliography}

\end{document}
